% !TEX TS-program = xelatex
% !TEX encoding = UTF-8 Unicode
% !Mode:: "TeX:UTF-8"

\documentclass{resume}
\usepackage{zh_CN-Adobefonts_external} % Simplified Chinese Support using external fonts (./fonts/zh_CN-Adobe/)
% \usepackage{NotoSansSC_external}
% \usepackage{NotoSerifCJKsc_external}
% \usepackage{zh_CN-Adobefonts_internal} % Simplified Chinese Support using system fonts
\usepackage{linespacing_fix} % disable extra space before next section
\usepackage{cite}

\begin{document}
\pagenumbering{gobble} % suppress displaying page number

\name{俞磊}

\basicInfo{
  \email{yulei@stu.hit.edu.cn} \textperiodcentered\
  \phone{15336376710} \textperiodcentered\
  \qq{1845385862}
  % \linkedin[billryan8]{https://www.linkedin.com/in/billryan8}
}

\section{\faGraduationCap\  教育背景}
\datedsubsection{\textbf{哈尔滨工业大学}}{2023 -- 至今}
\textit{在读硕士研究生}\ 网络空间安全, 预计 2026年毕业
\datedsubsection{\textbf{哈尔滨工业大学}}{2019 -- 2023}
\textit{学士}\ 计算机科学与技术

\section{\faUsers\ 实习经历}

\datedsubsection{\textbf{华为 数据通信产品线}}{2025年7月 -- 至今}
\keypoints{C, C++}
DHCP组件开发者测试框架扩充
\begin{itemize}
  \item 依据不同款型特性的部署需求，重写编译构建测试目标，检出1例主线版本中致命错误。
  \item 编写多进程款型及轻量化款型的测试入口及基础测试用例，代码覆盖率提升7\%。
  \item 增设门禁流水线LLT任务，部署各款型测试任务到日常代码合并请求行为。
  \item 编写DHCP报文Fuzz测试用例，对接内部测试平台，实现自动化Fuzz测试。
\end{itemize}

\datedsubsection{\textbf{华为 数据通信产品线}}{2024年4月 -- 2024年7月}
\keypoints{C, Lua}
Segment Routing路由协议组件迁移重构
\begin{itemize}
  \item 实现协议的按需下一跳组件的组建过程, 完成该组件从Ipv4到Ipv6的迁移
  \item 基于Netconf协议, 实现按需下一跳组件相关YANG协议节点
  \item 编写按需下一跳组件Lua脚本, 实现组件热更新
\end{itemize}

\section{\faUsers\ 项目经历}
\datedsubsection{\textbf{弹性云网调度系统}}{2023年4月 -- 2023年12月}
\keypoints{Java, Linux}
\begin{onehalfspacing}
多节点弹性自组织子网, 完成网络攻防任务
\begin{itemize}
  \item 负责设计弹性子网组建拆分过程以及通信协议
  \item 基于Netty完成网络通信功能开发
  \item 基于IPsec协议设计系统内可复用多跳匿名掩护通道
  \item 实现系统内节点根据任务需求, 自行征集不同节点组建子网单元执行任务
  \item 支持40节点以上的网络规模正常运作
\end{itemize}
\end{onehalfspacing}

\datedsubsection{\textbf{多层加密隧道组建系统}}{2022年9月 -- 2023年1月}
\keypoints{C++, Linux}
\begin{onehalfspacing}
Tor形式多层加密通道收集回流信息
\begin{itemize}
  \item 负责设计Tor通道上各节点与中心服务器通信协议
  \item 基于Muduo网络库, 使用RSA加密协议实现节点与中心服务器的通信与认证
  \item 支持负载均衡,可构建至少5节点的Tor通道
\end{itemize}
\end{onehalfspacing}

% \datedsubsection{\textbf{\LaTeX\ 简历模板}}{2015 年5月 -- 至今}
% \role{\LaTeX, Python}{个人项目}
% \begin{onehalfspacing}
% 优雅的 \LaTeX\ 简历模板, https://github.com/billryan/resume
% \begin{itemize}
%   \item 容易定制和扩展
%   \item 完善的 Unicode 字体支持，使用 \XeLaTeX\ 编译
%   \item 支持 FontAwesome 4.5.0
% \end{itemize}
% \end{onehalfspacing}

% Reference Test
%\datedsubsection{\textbf{Paper Title\cite{zaharia2012resilient}}}{May. 2015}
%An xxx optimized for xxx\cite{verma2015large}
%\begin{itemize}
%  \item main contribution
%\end{itemize}

% \section{\faCogs\ 专业技能}
% % increase linespacing [parsep=0.5ex]
% \begin{itemize}[parsep=0.5ex]
%   \item 编程语言: 熟练使用C++, 掌握Python, 了解 Java, Go
%   \item 平台: Linux
%   \item 开发工具: VsCode
% \end{itemize}

% \section{\faHeartO\ 获奖情况}
% \datedline{\textit{第一名}, xxx 比赛}{2013 年6 月}
% \datedline{其他奖项}{2015}

\section{\faInfo\ 自我评价}
% increase linespacing [parsep=0.5ex]
\begin{itemize}[parsep=0.5ex]
  \item 编程语言: 熟练使用C++, 掌握Python, 了解 Java, Lua, Go
  \item 熟悉Linux平台开发
  \item 熟悉TCP/IP网络协议与编程
  \item 对Key-value存储有一定了解,深入研究过leveldb实现原理
  \item 代码能力扎实,喜欢钻研和学习算法,\href{https://leetcode.cn/u/keshuigu/}{力扣Top2\%}
  \item 适应能力强,乐于团队协作与交流
\end{itemize}

%% Reference
%\newpage
%\bibliographystyle{IEEETran}
%\bibliography{mycite}
\end{document}
